\documentclass[11pt,a4paper]{article}

% ===== PACKAGES =====
\usepackage[utf8]{inputenc}
\usepackage[T1]{fontenc}
\usepackage{geometry}
\geometry{margin=2.5cm}
\usepackage{times}
\usepackage{graphicx}
\usepackage{booktabs}
\usepackage{array}
\usepackage{multirow}
\usepackage{enumitem}
\usepackage{xcolor}
\usepackage{hyperref}
\usepackage{tabularx}
\usepackage{caption}
\usepackage{float}
\usepackage{amsmath}
\usepackage{tcolorbox}
\usepackage{fancyhdr}

% ===== COLORS =====
\definecolor{darkBlue}{RGB}{25,60,120}
\definecolor{noteBlue}{RGB}{232,244,253}
\definecolor{noteBorder}{RGB}{66,133,244}
\definecolor{warnOrange}{RGB}{255,243,224}
\definecolor{warnBorder}{RGB}{230,126,34}
\definecolor{greenBox}{RGB}{232,250,237}
\definecolor{greenBorder}{RGB}{39,174,96}
\definecolor{purpleBox}{RGB}{240,232,253}
\definecolor{purpleBorder}{RGB}{142,68,173}

% ===== CUSTOM BOXES =====
\newtcolorbox{notebox}{
  colback=noteBlue, colframe=noteBorder,
  boxrule=0.5pt, left=6pt, right=6pt, top=4pt, bottom=4pt,
  fontupper=\small\itshape
}
\newtcolorbox{warnbox}{
  colback=warnOrange, colframe=warnBorder,
  boxrule=0.5pt, left=6pt, right=6pt, top=4pt, bottom=4pt,
  fontupper=\small\itshape
}
\newtcolorbox{greenbox}{
  colback=greenBox, colframe=greenBorder,
  boxrule=0.5pt, left=6pt, right=6pt, top=4pt, bottom=4pt,
  fontupper=\small\itshape
}
\newtcolorbox{purplebox}{
  colback=purpleBox, colframe=purpleBorder,
  boxrule=0.5pt, left=6pt, right=6pt, top=4pt, bottom=4pt,
  fontupper=\small\itshape
}

% ===== HEADER/FOOTER =====
\pagestyle{fancy}
\fancyhf{}
\fancyhead[C]{\small\textit{ProfsAgent --- Literature Gap Analysis \& System Vision}}
\fancyfoot[C]{\thepage}
\renewcommand{\headrulewidth}{0.4pt}

% ===== TITLE =====
\title{
\vspace{-1cm}
\textbf{ProfsAgent: Literature Gap Analysis}\\[0.3cm]
\large Identifying Opportunities for an Agentic AI System\\for Course Design and Delivery\\[0.5cm]
\normalsize Based on Sridhar et al.\ (2023) and Yao et al.\ (2026)
}

\author{
BTP Project --- ProfsAgent
}

\date{August 2026}

% ===== DOCUMENT =====
\begin{document}

\maketitle
\thispagestyle{fancy}

% ---------- ABSTRACT ----------
\begin{abstract}
Designing a new course involves defining learning outcomes, structuring topics, preparing lecture material, and developing assessments---a process that is time-consuming and particularly challenging for interdisciplinary or evolving subjects. Two recent papers explore using LLMs for this task: Sridhar et al.\ (2023) demonstrate prompt-engineered LO generation with GPT-4, and Yao et al.\ (2026) propose a multi-agent framework for end-to-end course material generation. This document identifies \textbf{13 key gaps} across both papers that inform the design of \textbf{ProfsAgent}---an agentic AI system that assists professors across the entire course-design lifecycle while keeping the professor in control of all academic decisions. The identified gaps directly translate into design requirements and differentiation points for ProfsAgent.
\end{abstract}

\noindent\rule{\textwidth}{0.4pt}
\tableofcontents
\noindent\rule{\textwidth}{0.4pt}
\vspace{0.3cm}

% ======================================================================
\section{Context: The ProfsAgent Vision}
% ======================================================================

\textbf{ProfsAgent} is an agentic AI system to assist professors across the course-design lifecycle. Depending on scope, the system may support:

\begin{itemize}[nosep]
    \item Curriculum planning and learning-outcome refinement
    \item Resource discovery and recommendation
    \item Lecture-note and slide generation
    \item Assignment and exam design
    \item Consistency checks across course components
\end{itemize}

The project may use \textbf{multiple specialized agents}, \textbf{retrieval-augmented generation (RAG)}, and \textbf{knowledge graphs}, while ensuring the professor remains in control. The two suggested readings serve as the foundational literature:

\begin{table}[H]
\centering
\caption{Overview of the two foundational papers.}
\begin{tabular}{@{}lp{5cm}p{5cm}@{}}
\toprule
\textbf{Dimension} & \textbf{Paper 1 (Sridhar, CMU)} & \textbf{Paper 2 (Yao, ASU)} \\
\midrule
Scope & LO generation only & Full course package \\
Architecture & Single model, single prompt & 6 agents, ADDIE pipeline \\
Models & GPT-4 only & gpt-4o, gpt-4o-mini, o1-preview \\
Courses tested & 1 (AI Practitioner) & 5 CS courses \\
Evaluators & 3 CS grad students + BERT & 5 expert instructors + LLM judges \\
Best quality & N/A (qualitative) & 3.74/5 (Co-Pilot mode) \\
Cost & Not reported & \$0.22--\$0.36 per course \\
\bottomrule
\end{tabular}
\end{table}

The following 13 gaps represent what these papers \textbf{did not do} or \textbf{did not do well}---and where ProfsAgent can differentiate itself.

% ======================================================================
\section{Gap 1: Limited Model Diversity and Benchmarking}
% ======================================================================

\textbf{Paper~1} uses only GPT-4. \textbf{Paper~2} tests three OpenAI models (gpt-4o, gpt-4o-mini, o1-preview) but excludes all non-OpenAI alternatives. Neither tests Claude, Gemini, Llama, Mistral, or any open-source model.

\subsection{Why This Matters for ProfsAgent}
Different models have different strengths: Claude for structured reasoning, Gemini for multi-modal tasks, open-source models for data privacy. ProfsAgent should be \textbf{model-agnostic}---allowing institutions to choose based on cost, privacy, and performance needs.

\subsection{Proposed Approach}
\begin{itemize}[nosep]
    \item Design ProfsAgent with a \textbf{pluggable model backend} (any OpenAI-compatible API).
    \item Benchmark the same pipeline across GPT-4o, Claude Sonnet~4, Gemini~2.5 Pro, and Llama~3.1.
    \item Explore \textbf{heterogeneous model assignment}---stronger models for complex agents (domain expert), cheaper models for formatting/TA tasks.
\end{itemize}

\begin{notebox}
\textbf{Key finding from Paper~2:} gpt-4o-mini matched gpt-4o at 1/16th the cost (Friedman test $p=0.789$). Cheaper models may suffice for many tasks.
\end{notebox}

% ======================================================================
\section{Gap 2: No Temperature or Hyperparameter Optimization}
% ======================================================================

\textbf{Paper~1} uses \texttt{temperature=0.7} (default) without justification. Neither paper experiments with hyperparameters.

\subsection{Proposed Approach}
\begin{itemize}[nosep]
    \item Conduct a temperature sweep (0.2, 0.5, 0.7, 0.9) measuring quality vs.\ diversity.
    \item Implement \textbf{Best-of-N}: generate $N$ candidate outputs, select the best via a quality scorer.
    \item Different agents/tasks may benefit from different temperatures (creative tasks $\to$ higher, structured tasks $\to$ lower).
\end{itemize}

% ======================================================================
\section{Gap 3: Region-Specific Validation Missing}
% ======================================================================

The BERT classifier in Paper~1 was trained on US university LOs. Educational conventions differ dramatically across countries.

\begin{table}[H]
\centering
\caption{LO conventions differ across educational systems.}
\begin{tabular}{@{}lp{3cm}p{3cm}p{3.5cm}@{}}
\toprule
\textbf{Aspect} & \textbf{US (ABET)} & \textbf{India (NBA/NAAC)} & \textbf{Europe (Bologna)} \\
\midrule
LO Format & ``Students will be able to\ldots'' & ``CO1: Design algorithms\ldots'' & ``Graduate can demonstrate\ldots'' \\
Taxonomy & Bloom's (6 levels) & Bloom's + SOLO & Dublin Descriptors / EQF \\
Accreditation & ABET & NBA / NAAC & EUR-ACE \\
\bottomrule
\end{tabular}
\end{table}

\subsection{Proposed Approach for ProfsAgent}
\begin{itemize}[nosep]
    \item Let the professor select their \textbf{country and accreditation body} during setup.
    \item Load region-specific LO templates, action verb lists, and validation rules.
    \item Use region-specific few-shot examples in prompts (not just CMU examples).
    \item \textbf{This is critical for Indian universities} (NBA/NAAC requirements are different from ABET).
\end{itemize}

% ======================================================================
\section{Gap 4: Homogeneous and Narrow Evaluator Panels}
% ======================================================================

\textbf{Paper~1} uses 3 CS grad students ($\kappa = 0.31$, fair). \textbf{Paper~2} uses 5 expert instructors (no $\kappa$ reported). Neither includes students.

\subsection{Proposed Evaluation Panel for ProfsAgent}

\begin{table}[H]
\centering
\caption{Multi-perspective annotator panel.}
\begin{tabular}{@{}lcl@{}}
\toprule
\textbf{Role} & \textbf{Count} & \textbf{Unique Perspective} \\
\midrule
Professor / Subject Expert & 1 & ``Is this pedagogically sound?'' \\
Strong Students (Top 20\%) & 3 & ``Is this challenging enough?'' \\
Average Students (Mid 40\%) & 3 & ``Do I understand what's expected?'' \\
\midrule
\textbf{Total} & \textbf{7} & \textbf{3 distinct perspectives} \\
\bottomrule
\end{tabular}
\end{table}

\begin{warnbox}
\textbf{Key Insight:} Average students are the most important validators. If course material isn't clear to them, it fails its primary purpose. Both papers completely lack this perspective.
\end{warnbox}

% ======================================================================
\section{Gap 5: No Baseline Comparison with Human-Created Materials}
% ======================================================================

Neither paper compares AI-generated outputs against \textbf{existing human-written materials} for the same course.

\subsection{Proposed Approach for ProfsAgent}
\begin{itemize}[nosep]
    \item Collect existing professor-written LOs, syllabi, and slides from partner courses.
    \item Conduct \textbf{blind evaluation}: reviewers score both AI and human outputs without knowing the source.
    \item Report head-to-head comparison---this is the most convincing evaluation for a publication.
\end{itemize}

% ======================================================================
\section{Gap 6: No Iterative Self-Critique in Generation}
% ======================================================================

\textbf{Paper~1} uses single-shot generation. \textbf{Paper~2}'s agents generate content but don't self-review before presenting to the professor.

\subsection{Proposed Approach for ProfsAgent}
Build a \textbf{Generate $\to$ Critique $\to$ Refine} loop into every agent:
\begin{enumerate}[nosep]
    \item \textbf{Generate:} Produce initial output (LOs, slides, assessments).
    \item \textbf{Self-Critique:} A reviewer agent evaluates the output against quality criteria, flagging issues (vague terms, Bloom's misalignment, multiple action verbs in one LO).
    \item \textbf{Refine:} Fix flagged issues automatically.
    \item \textbf{Present to Professor:} Only show the refined version, with flagged items highlighted for human review.
\end{enumerate}

This adds minimal cost (one extra LLM call) but consistently improves output quality.

% ======================================================================
\section{Gap 7: No Measurability Evaluation of Learning Objectives}
% ======================================================================

Both papers say LOs should be ``measurable'' but never evaluate this. Paper~1 explicitly lists it as future work.

\subsection{Why This Matters for ProfsAgent}
If an LO like \textit{``Understand neural networks''} is generated, there's no way to design an assessment for it. ProfsAgent should ensure every LO is assessment-mappable.

\subsection{Proposed Approach}
Build a \textbf{measurability checker} that scores each LO on:
\begin{itemize}[nosep]
    \item \textbf{Observable Behavior:} Can you watch a student perform this?
    \item \textbf{Assessment-Mappable:} Can you design a test question for this?
    \item \textbf{Success Criteria:} Is the degree of mastery specific?
\end{itemize}

\textbf{Example:}
\begin{itemize}[nosep]
    \item \textcolor{red}{\textbf{Fails:}} \textit{``Understand neural networks.''}
    \item \textcolor{greenBorder}{\textbf{Passes:}} \textit{``Explain the forward propagation process in a 3-layer neural network, identifying the role of weights, biases, and activation functions.''}
\end{itemize}

% ======================================================================
\section{Gap 8: No RAG / Grounding on Actual Course Content}
% ======================================================================

Both papers generate content from LLM parametric knowledge or module names alone. No access to textbooks, prior syllabi, lecture notes, or institutional data.

\subsection{Why This is Critical for ProfsAgent}
This is the \textbf{single biggest differentiator} ProfsAgent can offer. Professors have existing materials---textbooks they follow, previous syllabi, departmental guidelines. Ignoring these leads to generic, sometimes inaccurate content.

\subsection{Proposed Approach}
\begin{itemize}[nosep]
    \item Build a \textbf{RAG pipeline}: professors upload their textbook PDFs, prior syllabi, and institutional guidelines.
    \item Chunk and embed documents into a vector store (e.g., ChromaDB, Pinecone).
    \item Each generation agent retrieves relevant context before producing output.
    \item This directly solves the ``Python libraries'' vs.\ ``scikit-learn'' specificity problem from Paper~1.
\end{itemize}

\begin{greenbox}
\textbf{ProfsAgent advantage:} RAG-grounded generation produces institution-specific, textbook-aligned materials---not generic LLM output. This is what professors actually need.
\end{greenbox}

% ======================================================================
\section{Gap 9: CS/STEM Only --- No Cross-Discipline Testing}
% ======================================================================

\textbf{Paper~1}: 1 AI course. \textbf{Paper~2}: 5 CS courses. No humanities, natural sciences, arts, or professional disciplines.

\subsection{Why This Matters}
\begin{itemize}[nosep]
    \item \textbf{Humanities} (History, Philosophy): LOs emphasize argumentation, interpretation, critical thinking.
    \item \textbf{Natural Sciences} (Physics, Chemistry): LOs involve lab skills, experimental design.
    \item \textbf{Professional} (Law, Medicine): LOs involve case analysis, clinical reasoning.
\end{itemize}

ProfsAgent should demonstrate generalizability beyond CS. Even testing on 2--3 non-CS courses would significantly strengthen the contribution.

% ======================================================================
\section{Gap 10: No Visual or Multimedia Content Generation}
% ======================================================================

\textbf{Paper~2}'s slides are text-only LaTeX. Reviewers noted: \textit{``Slide design templates were overly uniform, limiting visual variety and student engagement.''}

\subsection{Why This Matters}
Real lecture slides contain diagrams, charts, code visualizations, and images. Text-only slides are pedagogically weaker.

\subsection{Proposed Approach for ProfsAgent}
\begin{itemize}[nosep]
    \item Generate slides as \textbf{PPTX} (via python-pptx) or \textbf{HTML/reveal.js}---not LaTeX.
    \item Integrate image generation for diagrams and illustrations.
    \item Use code execution for automated charts, plots, and visualizations.
    \item Use Mermaid for flowcharts and architecture diagrams.
    \item This eliminates LaTeX fragility (missing \texttt{[fragile]} tags, unescaped characters).
\end{itemize}

% ======================================================================
\section{Gap 11: No Real Classroom Deployment or Student Outcomes}
% ======================================================================

Neither paper deploys materials to real students. Both acknowledge this as critical future work.

\subsection{Proposed Approach for ProfsAgent}
While full deployment may be beyond BTP scope, ProfsAgent can:
\begin{itemize}[nosep]
    \item Pilot with \textbf{1--2 real courses} at the university---even partial deployment (e.g., AI-generated slides for 2 weeks of a course).
    \item Collect student feedback via surveys (clarity, engagement, usefulness).
    \item Compare with the professor's own materials for the same topics.
    \item This provides the strongest evidence for a publication at AIED or CHI.
\end{itemize}

% ======================================================================
\section{Gap 12: LLM-as-Judge is Unreliable for Educational Content}
% ======================================================================

\textbf{Paper~2} demonstrates that LLM evaluators assign tightly clustered scores (2.9--3.1) while humans produce discriminative scores (1.5--4.5). LLMs cannot distinguish good from mediocre instructional materials.

\subsection{Implication for ProfsAgent}
Do \textbf{not} rely on LLM-based quality scoring as the primary evaluation metric. Use:
\begin{itemize}[nosep]
    \item \textbf{Human evaluation} as the primary quality measure.
    \item \textbf{LLM reviewers} only for automated sanity checks (formatting, Bloom's verb detection, completeness).
    \item Consider fine-tuning a \textbf{domain-specific evaluation model} on expert-annotated educational materials.
\end{itemize}

% ======================================================================
\section{Gap 13: No Cross-Module Consistency or Curriculum Coherence}
% ======================================================================

Both papers generate content per-module \textbf{independently}. No verification that the full curriculum is coherent, non-overlapping, and properly sequenced.

\subsection{Proposed Approach for ProfsAgent}
Build a \textbf{Consistency Engine} that runs after all modules are generated:
\begin{itemize}[nosep]
    \item \textbf{Overlap detection:} Flag redundant LOs or content across modules.
    \item \textbf{Coverage matrix:} Modules $\times$ Bloom's Levels---identify gaps.
    \item \textbf{Prerequisite ordering:} Verify higher-level LOs appear after foundational ones.
    \item \textbf{Assessment alignment:} Every LO maps to at least one assessment item.
    \item Optionally build a \textbf{knowledge graph} of concepts and their dependencies to power this.
\end{itemize}

\begin{greenbox}
\textbf{This is a strong differentiator.} Neither paper attempts course-level coherence checking. A knowledge-graph-powered consistency engine would be a novel contribution.
\end{greenbox}

% ======================================================================
\section{Summary of Gaps and Relevance to ProfsAgent}
% ======================================================================

\begin{table}[H]
\centering
\caption{All 13 gaps with relevance to ProfsAgent system design.}
\small
\begin{tabular}{@{}clcp{4.5cm}@{}}
\toprule
\textbf{\#} & \textbf{Gap} & \textbf{Source} & \textbf{ProfsAgent Response} \\
\midrule
1 & Limited model diversity & P1+P2 & Pluggable model backend \\
2 & No hyperparameter optimization & P1 & Temperature tuning, Best-of-N \\
3 & Region-specific validation & P1 & Country/accreditation selector \\
4 & No student evaluators & P1+P2 & Multi-perspective panel (7 people) \\
5 & No human baseline comparison & P1+P2 & Blind AI vs.\ human evaluation \\
6 & No self-critique in generation & P1+P2 & Generate-Critique-Refine loop \\
7 & No measurability evaluation & P1 & Built-in measurability checker \\
8 & No RAG / content grounding & P1+P2 & RAG pipeline with professor's materials \\
9 & CS-only, no cross-discipline & P1+P2 & Test on 2--3 non-CS courses \\
10 & No visual/multimedia content & P2 & PPTX/HTML slides + image generation \\
11 & No classroom deployment & P1+P2 & Pilot with 1--2 real courses \\
12 & LLM-as-judge unreliable & P2 & Human-primary evaluation \\
13 & No curriculum coherence check & P1+P2 & Knowledge-graph consistency engine \\
\bottomrule
\end{tabular}
\end{table}

% ======================================================================
\section{ProfsAgent: High-Level Implementation Roadmap}
% ======================================================================

Based on the gaps identified and the BTP scope, we propose the following phased implementation plan:

\subsection{Phase 1: Foundation (Weeks 1--3)}
\begin{itemize}[nosep]
    \item Set up project repository, development environment, and CI/CD.
    \item Design the \textbf{system architecture}: agent framework, model abstraction layer, RAG pipeline.
    \item Build the \textbf{model abstraction layer}: pluggable backend supporting OpenAI, Claude, Gemini APIs.
    \item Define the \textbf{agent role specifications}: what each agent does, its system prompt, and inputs/outputs.
\end{itemize}

\subsection{Phase 2: Core Agent Pipeline (Weeks 4--7)}
\begin{itemize}[nosep]
    \item Implement \textbf{LO Generation Agent} with structured output (JSON schema for Behavior, Conditions, Degree, Bloom's level).
    \item Implement \textbf{Syllabus Generation Agent} that produces weekly topic plans from LOs.
    \item Implement \textbf{Slide Generation Agent} producing PPTX or HTML slides (not LaTeX).
    \item Implement \textbf{Assessment Generation Agent} for quizzes, assignments, and exam questions.
    \item Build the \textbf{Generate-Critique-Refine} loop into every agent.
\end{itemize}

\subsection{Phase 3: RAG \& Knowledge Pipeline (Weeks 6--9)}
\begin{itemize}[nosep]
    \item Build document ingestion pipeline (PDF $\to$ chunks $\to$ embeddings).
    \item Set up vector store (ChromaDB or FAISS).
    \item Integrate RAG retrieval into each agent's generation context.
    \item Optionally: build a \textbf{concept knowledge graph} from syllabus topics for prerequisite ordering.
\end{itemize}

\subsection{Phase 4: Consistency Engine (Weeks 8--10)}
\begin{itemize}[nosep]
    \item Build the \textbf{cross-module consistency checker}: overlap detection, Bloom's coverage matrix.
    \item Build the \textbf{LO-Assessment alignment verifier}: every LO maps to $\geq$1 assessment.
    \item Build the \textbf{prerequisite ordering checker}: verify course sequencing makes sense.
    \item Build the \textbf{measurability checker} for LOs.
\end{itemize}

\subsection{Phase 5: Professor Interface (Weeks 9--12)}
\begin{itemize}[nosep]
    \item Build a web-based UI (Streamlit, Gradio, or Next.js) for professor interaction.
    \item Support: course setup (name, goals, accreditation body), material review and editing, approval workflow.
    \item Implement the \textbf{Co-Pilot mode}: professor reviews and refines at each stage.
    \item Export generated materials as downloadable packages (PPTX, PDF, DOCX).
\end{itemize}

\subsection{Phase 6: Evaluation \& Publication (Weeks 11--14)}
\begin{itemize}[nosep]
    \item Recruit multi-perspective evaluation panel (1 prof + 3 strong + 3 average students).
    \item Run on 3--5 courses across CS and at least 1 non-CS discipline.
    \item Conduct blind AI vs.\ human baseline comparison.
    \item Report inter-rater reliability ($\kappa$), quality scores, cost analysis.
    \item Pilot deployment in 1 real course (even partially).
    \item Write paper targeting AIED or CHI.
\end{itemize}

\begin{purplebox}
\textbf{Note:} Phases overlap. RAG development (Phase~3) can run in parallel with core agent development (Phase~2). The UI (Phase~5) can start with a simple Streamlit prototype and be refined iteratively.
\end{purplebox}

% ======================================================================
% REFERENCES
% ======================================================================

\begin{thebibliography}{10}

\bibitem{sridhar2023}
P.~Sridhar, A.~Doyle, A.~Agarwal, C.~Bogart, J.~Savelka, and M.~Sakr.
\newblock Harnessing {LLMs} in curricular design: Using {GPT-4} to support authoring of learning objectives.
\newblock In \textit{Workshop on Empowering Education with LLMs, AIED}, 2023.
\newblock arXiv:2306.17459.

\bibitem{yao2026}
H.~Yao, W.~Xu, J.~Turnau, N.~Kellam, and H.~Wei.
\newblock Instructional Agents: {LLM} Agents on automated course material generation for teaching faculties.
\newblock Arizona State University, 2026.
\newblock arXiv:2508.19611.

\bibitem{li2022}
Y.~Li, M.~Rakovic, B.~X.~Poh, D.~Gasevic, and G.~Chen.
\newblock Automatic classification of learning objectives based on {Bloom's} taxonomy.
\newblock In \textit{Proceedings of EDM}, pp.~530--537, 2022.

\bibitem{bloom1956}
B.~S. Bloom, M.~D. Engelhart, E.~Furst, W.~H. Hill, and D.~R. Krathwohl.
\newblock \textit{Taxonomy of Educational Objectives: Handbook {I} --- Cognitive Domain}.
\newblock David McKay, New York, 1956.

\bibitem{krathwohl2002}
D.~R. Krathwohl.
\newblock A revision of {Bloom's} taxonomy: An overview.
\newblock \textit{Theory into Practice}, 41:212--218, 2002.

\bibitem{gagne2005}
R.~M. Gagne, W.~W. Wager, K.~C. Golas, J.~M. Keller, and J.~D. Russell.
\newblock \textit{Principles of Instructional Design}, 2005.

\bibitem{branch2009}
R.~M. Branch and {\.I}.~Varank.
\newblock \textit{Instructional Design: The ADDIE Approach}, volume 722.
\newblock Springer, 2009.

\bibitem{qm2023}
Quality Matters.
\newblock \textit{Higher Education Rubric, Seventh Edition}.
\newblock MarylandOnline, Inc., 2023.

\end{thebibliography}

\end{document}
